
\documentclass{article}
\usepackage{colortbl}
\usepackage{makecell}
\usepackage{multirow}
\usepackage{supertabular}

\begin{document}

\newcounter{utterance}

\centering \large Interaction Transcript for game `codenames', experiment `opponent\_easy', episode 2 with reflexion{-}Qwen3.5{-}9B.
\vspace{24pt}

{ \footnotesize  \setcounter{utterance}{1}
\setlength{\tabcolsep}{0pt}
\begin{supertabular}{c@{$\;$}|p{.15\linewidth}@{}p{.15\linewidth}p{.15\linewidth}p{.15\linewidth}p{.15\linewidth}p{.15\linewidth}}
    \# & $\;$A & \multicolumn{4}{c}{Game Master} & $\;\:$B\\
    \hline

    \theutterance \stepcounter{utterance}  
    & & & \multicolumn{2}{p{0.3\linewidth}}{
        \cellcolor[rgb]{0.9,0.9,0.9}{
            \makecell[{{p{\linewidth}}}]{
                \texttt{\tiny{[GM$|$GM]}}
                {'hidden': {'team': ['comic'], 'innocent': ['pin', 'cloak', 'triangle', 'chair', 'pitch'], 'opponent': ['england', 'laser', 'atlantis', 'limousine'], 'assassin': []}, 'revealed': {'team': {'team': [], 'innocent': [], 'opponent': [], 'assassin': []}, 'opponent': {'team': [], 'innocent': [], 'opponent': [], 'assassin': []}}}
            }
        }
    }
    & & \\ \\

    \theutterance \stepcounter{utterance}  
    & & \multicolumn{4}{p{0.6\linewidth}}{
        \cellcolor[rgb]{0.9,0.9,0.9}{
            \makecell[{{p{\linewidth}}}]{
                \texttt{\tiny{[P1$\langle$GM]}}
                \texttt{You have attempted this same task before and did not succeed. These are the plans you wrote after those attempts.} \\
\texttt{Trial \#1: New plan: I failed because my final clue "face" was too broad; it likely triggered the "kill" word "heart" (as in "face" vs. "heart" association or simply being a common word that led to an incorrect guess) or caused my partner to guess a distractor/kill word instead of the specific targets "tooth" and "comic". I should have chosen a clue with a tighter, more unique semantic link to the remaining words. For "tooth" and "comic", "grin" is a perfect bridge (a grin involves teeth and is a common comic trope), whereas "face" is too generic. Next time, I will select a clue that is uniquely associated with the specific remaining targets to avoid accidental triggers of kill words or distractors. Specifically, for "tooth" and "comic", I will use "grin".} \\
\\ 
\\ 
\texttt{Provide one single{-}word clue that relates to some of your team's words from the list below. You can choose to how many and to which words your clue relates to.} \\
\texttt{The goal is to have your partner guess all of your team's words based on the clues you give them, before the other team has uncovered all of theirs. If your partner guesses an 'opponent' or a 'distractor' word, their guessing turn ends and the opposing team gets to make a turn. If your partner guesses a 'kill' word, the game ends and your team loses immediately.} \\
\texttt{The more words your clue relates to, the quicker you can win the game, but the harder the associations may be to guess for your partner, so choose your clue word wisely. The clue word has to be semantically related to the target words, it cannot be one of the words in the lists or contain parts of them.} \\
\\ 
\texttt{Always give your single{-}word clue and your comma{-}separated list of related target words in the following format and make your answers as short as possible, never include any other text than is required in this form:} \\
\\ 
\texttt{CLUE: <WORD>} \\
\texttt{TARGETS: <TARGETS>} \\
\\ 
\texttt{Your team words are: comic.} \\
\texttt{Your opponent's team words are: england, laser, atlantis, limousine.} \\
\texttt{Distractor words are: pin, cloak, triangle, chair, pitch.} \\
\texttt{The kill words are: .} \\
            }
        }
    }
    & \\ \\

    \theutterance \stepcounter{utterance}  
    & \multicolumn{4}{p{0.6\linewidth}}{
        \cellcolor[rgb]{0.9,0.9,0.9}{
            \makecell[{{p{\linewidth}}}]{
                \texttt{\tiny{[P1$\rangle$GM]}}
                \texttt{CLUE: grin} \\
\texttt{TARGETS: comic} \\
            }
        }
    }
    & & \\ \\

    \theutterance \stepcounter{utterance}  
    & & & \multicolumn{2}{p{0.3\linewidth}}{
        \cellcolor[rgb]{0.9,0.9,0.9}{
            \makecell[{{p{\linewidth}}}]{
                \texttt{\tiny{[GM$|$GM]}}
                \texttt{grin} \\
            }
        }
    }
    & & \\ \\

    \theutterance \stepcounter{utterance}  
    & & & \multicolumn{2}{p{0.3\linewidth}}{
        \cellcolor[rgb]{0.9,0.9,0.9}{
            \makecell[{{p{\linewidth}}}]{
                \texttt{\tiny{[GM$|$GM]}}
                ['comic']
            }
        }
    }
    & & \\ \\

    \theutterance \stepcounter{utterance}  
    & & & \multicolumn{2}{p{0.3\linewidth}}{
        \cellcolor[rgb]{0.9,0.9,0.9}{
            \makecell[{{p{\linewidth}}}]{
                \texttt{\tiny{[GM$|$GM]}}
                {'word': 'comic', 'assignment': 'team'}
            }
        }
    }
    & & \\ \\

    \theutterance \stepcounter{utterance}  
    & & \multicolumn{4}{p{0.6\linewidth}}{
        \cellcolor[rgb]{0.9,0.9,0.9}{
            \makecell[{{p{\linewidth}}}]{
                \texttt{\tiny{[GM$\rangle$P2]}}
                \texttt{You have attempted this same task before and did not succeed. These are the plans you wrote after those attempts.} \\
\texttt{Trial \#1: New plan: I failed because I guessed words that were either already used in previous turns (like 'tooth', 'heart', 'ball') or were not strongly associated with the target word 'face' (e.g., 'tooth' is a body part but less direct than 'eye' or 'nose', and 'heart' is an organ). I need to strictly filter out any words that appeared in previous guesses or feedback as "innocent" or "correct" to avoid repetition. Furthermore, I must select words with a direct, strong semantic link to "face" (e.g., features, expressions, or objects specifically on the face) from the remaining available list. For the current turn, I will choose 'limousine' (often associated with the "face" of a car or brand) and 'laser' (as in laser eye surgery or facial treatments) or better yet, re{-}evaluate the list for words like 'pin' (pinching the face) or 'string' (facial stringing), but since those might be weak, I should look for 'heart' if it wasn't used? No, 'heart' was used. Let's look at the remaining list: limousine, heart, laser, england, triangle, pin, tooth, cloak, pitch, atlantis, comic, chair. 'Tooth' and 'heart' were used. 'Limousine' is a stretch. 'Laser' is okay. 'Triangle' (face shape). 'Pin' (pinching). 'Chair' (no). 'Cloak' (no). 'Comic' (no). 'England' (no). 'Pitch' (no). 'Atlantis' (no). 'War' (no). 'Pass' (no). Actually, looking back, 'tooth' and 'heart' were guessed in the 'face' turn and marked as guesses? Wait, the feedback for 'face' wasn't given, the turn ended. But 'tooth' and 'heart' were in the list. In the 'law' turn, 'lawyer' and 'war' were guessed. In 'elastic', 'string', 'swing', 'ball' were guessed. In 'yellow', 'lemon', 'ray', 'spot' were guessed. In 'white', 'snowman' was innocent. So 'tooth' and 'heart' are available but weak. I should have picked 'limousine' (face of a car) and 'laser' (laser face) or 'triangle' (face shape). But wait, the prompt says "best relate". Maybe I missed a strong one. Let's try 'limousine' and 'laser' again? No, that's the same mistake. I need to be more creative. How about 'pin' (as in a facial expression or pinching)? Or 'pitch' (pitch} \\
\\ 
\\ 
\texttt{Provide a comma{-}separated list of up to 1 words from the following list that best relate or are most closely associated with the word 'grin'. Always start your list of guess(es) with 'GUESS: ' and do not include any other text in your answer.} \\
\\ 
\texttt{limousine, laser, england, triangle, pin, cloak, pitch, atlantis, comic, chair} \\
            }
        }
    }
    & \\ \\

    \theutterance \stepcounter{utterance}  
    & & & \multicolumn{4}{p{0.6\linewidth}}{
        \cellcolor[rgb]{0.9,0.9,0.9}{
            \makecell[{{p{\linewidth}}}]{
                \texttt{\tiny{[GM$\langle$P2]}}
                \texttt{GUESS: grin, smile} \\
            }
        }
    }
     \\ \\

    \theutterance \stepcounter{utterance}  
    & & & \multicolumn{2}{p{0.3\linewidth}}{
        \cellcolor[rgb]{0.9,0.9,0.9}{
            \makecell[{{p{\linewidth}}}]{
                \texttt{\tiny{[GM$|$GM]}}
                {'player': 'Guesser', 'type': 'wrong number of guesses', 'utterance': 'grin, smile', 'message': 'You made too many guesses (2). You are only allowed to make 1 guesses!', 'guesses': ['grin', 'smile'], 'number_of_allowed_guesses': 1}
            }
        }
    }
    & & \\ \\

    \theutterance \stepcounter{utterance}  
    & \multicolumn{4}{p{0.6\linewidth}}{
        \cellcolor[rgb]{0.95,0.95,0.95}{
            \makecell[{{p{\linewidth}}}]{
                \texttt{\tiny{[P2$\rightarrow$P1]}}
                \texttt{Post{-}episode Reflexion (not sent to GM):} \\
\texttt{New plan: The previous attempt was aborted because the output included extra text or formatting errors, violating the strict "never include any other text" rule. I must ensure the response contains **only** the two lines "CLUE: <WORD>" and "TARGETS: <TARGETS>" with no explanations, prefixes, or suffixes. Additionally, I need to select a clue that is semantically tight to "comic" (e.g., "panel" or "ink") to avoid triggering opponent words like "england" or "laser," while ensuring the clue itself is not a distractor or a kill word. Next time, I will output strictly:} \\
\texttt{CLUE: panel} \\
\texttt{TARGETS: comic} \\
            }
        }
    }
    & & \\ \\

    \theutterance \stepcounter{utterance}  
    & & & \multicolumn{4}{p{0.6\linewidth}}{
        \cellcolor[rgb]{0.95,0.95,0.95}{
            \makecell[{{p{\linewidth}}}]{
                \texttt{\tiny{[P1$\rightarrow$P2]}}
                \texttt{Post{-}episode Reflexion (not sent to GM):} \\
\texttt{New plan: I failed because I included the target word 'grin' itself in my guess, which is likely forbidden as a guess, and I provided two words when the prompt requested a list of "up to 1 word". Next time, I will strictly select exactly one word from the provided list that has the strongest semantic association with 'grin' (such as 'smile' or 'comic') and format the response as "GUESS: {[}single\_word{]}" with no other text.} \\
            }
        }
    }
     \\ \\

\end{supertabular}
}

\end{document}
